\chapter{Leukopenia (Low White Blood Cell Count)}

\section*{Definition}
Leukopenia is a lower-than-normal total white blood cell count. Its significance depends on which cell line is reduced, the absolute neutrophil count, the trend, medications, and the presence of infection or other blood-count abnormalities.

\section*{What Happens in Your Body}
White blood cells coordinate innate and adaptive immune defense. Reduced numbers may impair protection against infection, while an isolated mild reduction can be transient or clinically insignificant. Neutropenia carries the greatest immediate bacterial and fungal infection risk.

\section*{Causes}
\begin{itemize}
  \item Viral infection, chemotherapy, radiation, marrow disorders, leukemia, lymphoma, and marrow infiltration.
  \item Medicines such as antithyroid drugs, clozapine, anticonvulsants, antibiotics, and immunosuppressants.
  \item Autoimmune disease, hypersplenism, severe nutritional deficiency, and congenital disorders.
  \item Laboratory variation or a benign constitutional low count in some individuals.
\end{itemize}

\section*{When to See a Doctor}
Seek urgent care for fever, chills, confusion, shortness of breath, severe sore throat, painful urination, abdominal pain, or rapidly worsening illness, particularly when the ANC is low. Persistent or progressively falling counts require medical evaluation.

\section*{Related Symptoms}
\begin{itemize}
  \item Leukopenia may be asymptomatic; recurrent infections, mouth ulcers, fatigue, fever, or bruising may prompt testing.
  \item Anemia or low platelets occurring with leukopenia raises concern for a broader marrow process.
\end{itemize}

\section*{Self-Care / Home Management}
Do not stop prescribed medicine without advice. Use careful hand hygiene, avoid known infectious exposures, and obtain prompt assessment for fever. A white-cell count alone does not determine infection risk; ask for the ANC and the clinical interpretation.

\section*{How Your Doctor Will Evaluate You}
\begin{itemize}
  \item Repeat complete blood count with differential and compare prior results; review medications, recent infections, nutrition, travel, and family history.
  \item Examine for lymph-node or spleen enlargement and signs of infection; test for nutritional, viral, autoimmune, or endocrine causes when indicated.
  \item Hematology referral and bone-marrow testing may be needed for persistent, severe, unexplained, or multilineage abnormalities.
\end{itemize}

\section*{Treatment Options}
Treatment is directed at the cause and may include stopping a causative drug under supervision, treating infection or nutritional deficiency, using growth-factor therapy in selected cases, or treating an underlying marrow or autoimmune disorder. Fever with severe neutropenia requires emergency antimicrobial treatment.

\section*{References}
\begin{thebibliography}{99}
\bibitem{leukopenia:ref1} Newburger PE, Dale DC. Evaluation and management of patients with isolated neutropenia. \textit{Semin Hematol}. 2013;50:198--206.
\bibitem{leukopenia:ref2} Tefferi A, et al. Approach to the adult with unexplained cytopenia. \textit{Mayo Clin Proc}. 2023;98: appropriate clinical review.
\end{thebibliography}
